\documentclass[11pt,letterpaper]{article}
\usepackage[margin=0.83in]{geometry}
\usepackage{fontspec,unicode-math}
\setmainfont{Latin Modern Roman}
\setsansfont{Latin Modern Sans}
\setmonofont{DejaVu Sans Mono}[Scale=0.80]
\setmathfont{Latin Modern Math}
\usepackage{microtype,amsmath,amsthm,booktabs,longtable,tabularx,array}
\usepackage[dvipsnames]{xcolor}
\usepackage{enumitem,fvextra,titlesec,fancyhdr}
\usepackage[unicode,breaklinks,colorlinks,linkcolor=MidnightBlue,urlcolor=MidnightBlue]{hyperref}
\usepackage{xurl}
\definecolor{ink}{HTML}{18344D}
\definecolor{accent}{HTML}{137C83}
\titleformat{\section}{\Large\sffamily\bfseries\color{ink}}{\thesection}{0.6em}{}
\titleformat{\subsection}{\large\sffamily\bfseries\color{ink}}{\thesubsection}{0.6em}{}
\setlength{\parindent}{0pt}
\setlength{\parskip}{5.5pt}
\setlist{nosep,leftmargin=1.4em,topsep=5pt}
\setlength{\emergencystretch}{3em}
\pagestyle{fancy}
\fancyhf{}
\fancyhead[L]{\small\sffamily Additional Lean proof simplifications}
\fancyhead[R]{\small\sffamily ProveIt / 26 October 2026}
\fancyfoot[C]{\small\thepage}
\setlength{\headheight}{14pt}
\newcommand{\code}[1]{\texttt{\detokenize{#1}}}
\newcommand{\repo}{https://github.com/VladimirReshetnikov/ProveIt/blob/7fa06d530df54d4aeb2e18027409df082234c69e}
\newcommand{\src}[1]{\href{\repo/Computability/DiophantineConsistency/Lean/DiophantineConsistency/#1.lean}{\code{#1.lean}}}
\newcommand{\ml}[1]{\href{https://github.com/leanprover-community/mathlib4/blob/81a5d257c8e410db227a6665ed08f64fea08e997/Mathlib/#1.lean}{mathlib source}}
\DefineVerbatimEnvironment{Lean}{Verbatim}{fontsize=\small,breaklines=true,breakanywhere=true,frame=single,rulecolor=\color{black!15},framesep=7pt}
\newtheorem{proposition}{Proposition}
\newtheorem{lemma}{Lemma}
\hypersetup{pdftitle={Additional Lean Proof Simplifications for ProveIt},
pdfauthor={Research and verification report prepared for Vladimir Reshetnikov},
pdfsubject={Checked proposals beyond the incoming proof-simplification packets}}
\begin{document}
\begin{titlepage}
\vspace*{0.8in}
{\color{accent}\sffamily\bfseries PROVEIT / DIOPHANTINE CONSISTENCY\par}
\vspace{0.35in}
{\Huge\sffamily\bfseries\color{ink}Additional Lean\\[7pt]Proof Simplifications\par}
\vspace{0.3in}
{\Large Checked proposals beyond the incoming packets\par}
\vspace{0.3in}
{\large Prepared for Vladimir Reshetnikov\\26 October 2026\par}
\vspace{0.4in}
\hrule
\vspace{0.2in}
\textbf{Main result.} A proposed patch changes nine existing modules and adds
one shared module. It saves \textbf{106 physical lines}, counting every added
helper, wrapper, comment, and blank line. Its main mathematical improvements
preserve signed rational roots during descent, reuse fundamental-unit
classification for small Pell ordinates, and compose sign-ambiguous residues.
Its representation improvements consume existing certificate equivalences
without changing the represented data.

\textbf{Novelty scope.} The comparison covers the seven ZIP archives present in
\code{docs/incoming}, plus the earlier packet referenced by that directory's
README, at the exact repository snapshot below. This is a claim of additional
repository proposals relative to that inspected material, not a claim of
priority for the underlying mathematics or Lean techniques.

\textbf{Reproducibility.} The accompanying archive contains the LaTeX and PDF,
the patch and complete proposed modules, original-versus-proposed metrics,
nine complete-value equality checks for representation constructors, Lean
validation logs, source and archive hashes, a detailed exclusion audit, and
instructions for replay. An additional nine-line permutation proposal is kept
separate because its full repository client was not compiled.
\vfill
{\small
Repository commit:\par
\code{7fa06d530df54d4aeb2e18027409df082234c69e}\par
Lean: \code{4.32.0}\quad
mathlib: \code{81a5d257c8e410db227a6665ed08f64fea08e997}\par}
\end{titlepage}
\setcounter{tocdepth}{1}
\tableofcontents
\newpage

\section{Scope, evidence, and results}

The request is to shorten and simplify existing proofs in
\href{https://github.com/VladimirReshetnikov/ProveIt/tree/7fa06d530df54d4aeb2e18027409df082234c69e/Computability/DiophantineConsistency}{\code{Computability/DiophantineConsistency}},
while avoiding proposals already present in
\href{https://github.com/VladimirReshetnikov/ProveIt/tree/7fa06d530df54d4aeb2e18027409df082234c69e/docs/incoming}{\code{docs/incoming}}.
All source locations and comparisons in this report refer to commit
\code{7fa06d530df54d4aeb2e18027409df082234c69e}.
The live branch can subsequently acquire further proposals.

The main library directory contains 391 Lean files and 71,601 physical lines
in the inspected snapshot. Including its entry point and checking directory,
the broader Lean directory contains 574 Lean files. These counts establish
the scope of the source survey; they do not mean every theorem was individually
reproved or every module rebuilt.

The incoming audit covers seven archives containing 1,990 members. It identifies
177 report/proposal artifacts and separately reviews the 60-file placed packet
referenced by the incoming README. Bundled source snapshots were used to
understand earlier recommendations, but were not mistaken for recommendations
themselves. The resulting exclusion map has 57 families, with archive names,
member paths, section locations, and prior claimed validation status.
The complete map is supplied as \code{audit/novelty_audit.md} and JSON.

\subsection{What was deliberately excluded}

Earlier packets already propose core list-lookup reuse, recursor congruence,
assignment transport and witness-block splitting, generic sum-of-squares
lemmas, shared degree and cost proofs, scanner simulation, circuit power
certificates, coefficient envelopes, partial-decoder injectivity, structured
source facts, modular cast cleanup, and several finite-tuple interfaces.
Those proposals are not presented again as contributions.

Some techniques below use general library facts that earlier reports also
mention. The contribution must therefore be identified more precisely than
using \code{simp} or factoring a helper. For example, the pairing proposal
uses an inverse coding relation to remove particular existential coordinates.
It does not repeat the earlier recommendation to split the witness block of
\code{Rep.conjunction}. Likewise, the Boolean proposal consumes an existing
quadratic zero-certificate equivalence; it does not redesign a scanner's
numeric Boolean interface.

No matching proposal was found in the inspected reports, notes, candidate
code, or placed packet for the seven directions in this report. This is a
bounded audit result. It does not rule out a similar idea elsewhere in the
repository's history, in material not present at this snapshot, or in the
wider Lean community.

\subsection{Measured source changes}

The main patch keeps all pre-existing declaration signatures in the changed
files. Its representation changes additionally have complete-value equality
checks, discussed in Section~\ref{sec:validation}.

\begin{center}
\small
\begin{tabularx}{\linewidth}{@{}Xrrr@{}}
\toprule
Module or group & Before & After & Saved\\
\midrule
NearSquareFinalizer & 255 & 251 & 4\\
MarginFinalizer & 177 & 135 & 42\\
PairingRepresentation & 88 & 94 & $-6$\\
QuadraticNatPairing & 158 & 139 & 19\\
QuadraticTreeCodeGraphs & 150 & 136 & 14\\
QuadraticBoolean & 136 & 126 & 10\\
NativePellBootstrap & 685 & 677 & 8\\
NativePellStepDown & 248 & 196 & 52\\
SignedPellResidue & 102 & 96 & 6\\
SignedModEq (new) & 0 & 43 & $-43$\\
\midrule
\textbf{Main patch} & \textbf{1,999} & \textbf{1,893} & \textbf{106}\\
\midrule
WangTiles942Permutation (optional) & 264 & 255 & 9\\
\bottomrule
\end{tabularx}
\end{center}

These are physical source lines, including comments and blank lines, computed
from the actual delivered files. Negative entries charge shared infrastructure
to the change that introduces it. The 106-line result is a 5.3\% reduction in
the selected production source, not a 5.3\% reduction in the entire library.
The optional nine lines are not included in the main total.

Source length is one measure. The larger improvement is that several proofs
now expose the mathematical reason for the conclusion: a denominator equal
to one, growth of Pell powers, transitivity of a signed residue relation, or
an already-proved certificate equivalence. No compilation-speed or
proof-term-size improvement is claimed.

\section{Preserve the exact rational root during integral descent}
\label{sec:roots}

\subsection{The avoidable loss of information}

Both \src{NearSquareFinalizer} and \src{MarginFinalizer} reconstruct rational
numbers \(r,s\) satisfying
\[
 r^2=a,\qquad s^2=b,\qquad
 \frac{\theta}{\omega}=r+w s.
\]
The existing proofs then use \code{Rat.isSquare_intCast_iff} to obtain
\emph{some} integer roots \(r_0,s_0\). They must recover the relationship between
the roots in the linear equation and those new witnesses:
\[
r=\pm r_0,\qquad s=\pm s_0.
\]
This leads to four combinations, each carrying essentially the same final
cast and multiplication argument. The original evidence was more informative
than the intermediate \code{IsSquare} statement.

The replacement retains each rational witness. A rational number whose
nonzero natural power is an integer already has denominator one. Its numerator
is therefore the exact signed integer needed in the linear equation.

\begin{lemma}
Let \(q\in\mathbb Q\), \(z\in\mathbb Z\), and \(n\in\mathbb N\setminus\{0\}\).
If \(q^n=z\), then \(q=\operatorname{num}(q)\), with the right side embedded in
\(\mathbb Q\).
\end{lemma}

\begin{proof}
Write \(q\) in reduced form. The rational denominator API gives
\[
 \operatorname{den}(q^n)=\operatorname{den}(q)^n.
\]
An embedded integer has denominator one. Hence
\(\operatorname{den}(q)^n=1\), which forces
\(\operatorname{den}(q)=1\) because \(n\ne0\). The numerator identity follows.
This reasoning preserves the sign of \(q\).
\end{proof}

The checked implementation is short:
\begin{Lean}
theorem rat_intCast_num_of_pow {q : ℚ} {z : ℤ} {n : ℕ}
    (hn : n ≠ 0) (h : q ^ n = (z : ℚ)) :
    (q.num : ℚ) = q := by
  apply Rat.coe_int_num_of_den_eq_one
  simpa [hn] using congrArg Rat.den h
\end{Lean}

The proof uses the denominator identities for powers and embedded integers,
the natural-power equation for one, and
\code{Rat.coe_int_num_of_den_eq_one}; see
\ml{Data/Rat/Defs} and the Lean rational-number definitions included in the
pinned toolchain. The hypothesis \(n\ne0\) is essential: every rational number
has zeroth power one.

\subsection{The client proof keeps both signs automatically}

After the existing rational reconstruction, the new finalizer proof uses:
\begin{Lean}
have hrInt := rat_intCast_num_of_pow (by decide : 2 ≠ 0) hr
have hsInt := rat_intCast_num_of_pow (by decide : 2 ≠ 0) hs
rw [← hrInt] at hr hsum
rw [← hsInt] at hs hsum
refine ⟨r.num, s.num, ?_, ?_, ?_⟩
· exact_mod_cast hr
· exact_mod_cast hs
· exact_mod_cast
    ((div_eq_iff hoQ).mp hsum).trans (mul_comm _ _)
\end{Lean}

The integer witnesses are the numerators of the very roots that already
satisfy the sum equation. There is no sign search. Negative roots, zero
radicands where otherwise permitted, and either sign of the nonzero weight
remain covered.

This is more specific than a general cast cleanup. It replaces an intermediate
existence assertion about an integer square root by integrality of the
particular rational root. The latter is exactly the information the subsequent
linear relation needs.

\subsection{Put the general separated theorem below both clients}

The two files also contain near-duplicate root reconstruction proofs.
Their actual common hypotheses are
\[
 \omega\ne0,\qquad w\ne0,\qquad a\ne w^2b,
 \qquad \operatorname{norm}(\theta,\omega,a,w,b)=0.
\]
These suffice to reconstruct integer \(r,s\) with
\[
 r^2=a,\qquad s^2=b,\qquad \theta=\omega(r+ws).
\]

The patch places this separated theorem in \code{NearSquareFinalizer},
where the norm and rational reconstruction already live.
\code{MarginFinalizer.integer_roots_of_separated_norm_zero} keeps its original
public name and signature as a direct wrapper.

The earlier bounded form in \code{NearSquareFinalizer} now obtains separation
from its existing hypotheses \(w>0\), \(a<w^2\), \(b>0\). Since \(b\) is an
integer, \(b\ge1\), so \(w^2b\ge w^2>a\). Its proof becomes an application of the
general theorem plus this elementary inequality.

This relocation follows the actual dependency direction. It adds no import
to either file and no cycle. It also preserves both original APIs.
The norm and final-polynomial definitions are unchanged, and the complete
revised files compile together. The net reduction is \textbf{46 lines},
including the power lemma, generalized theorem, specialization, and wrapper.

\section{A uniform small-ordinate criterion for Pell minimality}
\label{sec:pell}

\subsection{A stronger lemma replaces a discriminant-specific argument}

In \src{NativePellBootstrap}, \code{firstUnit_fundamental} concerns
\[
 d=Q(Q+1), \qquad (x,y)=(2Q+1,2), \qquad Q>0.
\]
The existing proof analyzes a competing solution, bounds its squared
ordinate, splits its ordinate into \(-1,0,1\), and excludes those possibilities
using adjacent-square inequalities. That elementary proof is valid. The
following general fact provides a shorter route with existing mathlib APIs.

\begin{proposition}
For any integer \(d\), let \(a=(x,y)\) solve \(x^2-dy^2=1\).
If \(x>1\) and \(0<y\le3\), then \(a\) is fundamental.
The bound three cannot be increased to four.
\end{proposition}

\begin{proof}
The existence of a solution with \(x>1\) implies that \(d\) is positive and
nonsquare. Let \(f=(u,v)\) be the fundamental positive solution.
The classification theorem expresses \(a=f^n\) for a natural \(n\).
Positive \(y\) excludes \(n=0\).

If \(n\ge2\), monotonicity of the ordinate along powers gives
\[
 a.y\ge (f^2).y=2uv.
\]
Fundamentality gives the integer bounds \(u\ge2\) and \(v\ge1\), hence
\(a.y\ge4\), a contradiction. Thus \(n=1\).

For sharpness, \(d=3\) has the solution \((7,4)\). It is not fundamental
because \((2,1)\) is a positive solution with a smaller first coordinate.
\end{proof}

The checked general theorem is:
\begin{Lean}
theorem isFundamental_of_y_le_three
    {d : ℤ} {a : Pell.Solution₁ d}
    (hx : 1 < a.x) (hy : 0 < a.y) (hsmall : a.y ≤ 3) :
    Pell.IsFundamental a := by
  obtain ⟨f, hf⟩ := Pell.IsFundamental.exists_of_not_isSquare
    (Pell.Solution₁.d_pos_of_one_lt_x hx)
    (Pell.Solution₁.d_nonsquare_of_one_lt_x hx)
  obtain ⟨n, rfl⟩ := hf.eq_pow_of_nonneg (by omega) hy.le
  have hn : n < 2 := by
    by_contra hn
    have h2 := hf.y_strictMono.monotone
      (show (2 : ℤ) ≤ n by omega)
    simp only [zpow_natCast, zpow_ofNat, pow_two,
      Pell.Solution₁.y_mul] at h2
    have hfx := hf.1
    have hfy := hf.2.1
    nlinarith only [h2, hsmall, hfx, hfy,
      mul_nonneg (show (0 : ℤ) ≤ f.x - 2 by omega)
        (show (0 : ℤ) ≤ f.y - 1 by omega)]
  interval_cases n
  · simp at hy
  · simpa using hf
\end{Lean}

The exact source file formats two of these expressions on longer lines;
the displayed version wraps them for readability. The source metrics count
the delivered file, not this presentation.

\subsection{Benefits and tradeoff}

The original first-unit theorem becomes a two-line proof body invoking the
criterion. Its nonblank proof body decreases from 30 lines to 2; charging
the new helper's 14 nonblank proof-body lines leaves a reduction of 14 by
that measure. Counting the complete production file, including documentation
and all declarations, gives \textbf{685 to 677 lines}, a saving of eight.

The new theorem is reusable for any positive Pell solution with ordinate
one, two, or three. Its sharpness witness is also checked in the standalone
artifact. This is a mathematical strengthening of the reusable interface,
rather than moving the old case analysis into a differently named lemma.

The tradeoff is explicit: the new proof uses mathlib's general existence and
classification of fundamental units, where the old proof establishes the
particular minimality directly. These theorems are already available through
the original imports. The patch adds no axiom or import, but it does change
which library results the proof uses. A project wishing to retain a visibly
elementary minimality proof may prefer the original. The primary source for
the reused results is \ml{NumberTheory/Pell}.

\section{Reuse the existing odd-index theorem}
\label{sec:odd}

There is a direct mathematical duplicate across
\src{NativePellStepDown} and \src{SignedPellResidue}.
The former proves that the auxiliary Pell index \(s\) is odd by assuming it
even, calculating a residue \((-1)^h\), proving that the base divides one,
and contradicting the lower bound on that base.

The latter already contains the precise reusable statement:
\[
 A>1,\quad A\mid x_A(n)\quad\Longrightarrow\quad n\text{ is odd},
\]
as \code{SignedPellResidue.odd_index_of_base_dvd}.
After auxiliary classification the step-down proof has
\(Iu=x_I(s)\), so the needed divisibility is immediate. The replacement is:

\begin{Lean}
have hsodd : Odd s :=
  SignedPellResidue.odd_index_of_base_dvd hI1
    (by rw [← hIu]; exact dvd_mul_right I u)
\end{Lean}

This reuses an existing result with exactly matching hypotheses. It needs no
general theorem search tactic and no reformulation of the public step-down
statement. The old local \code{hs0} proof of \(s>0\) is not used later and is
also removed. Its former input \code{hy} remains in the public signature,
so existing callers retain the same interface.

The additional import is acyclic:
\[
 \texttt{NativePellStepDown}\longrightarrow
 \texttt{SignedPellResidue}\longrightarrow
 \texttt{Diophantine.Common.PellInt}.
\]
In the combined patch the middle module also imports the new signed-congruence
module described next. No reverse dependency is introduced.

The novelty distinction is narrow but useful. Earlier packets discuss casting
whole Pell disjunctions and exposing source-register facts. This proposal
instead removes a duplicate number-theoretic argument by using an existing
odd-index theorem from a different file. It is an immediate adoption candidate
even if the broader sign interface is deferred.

\section{Compose congruence up to sign as a relation}
\label{sec:sign}

\subsection{The semantic operation is transitivity}

The last part of \code{NativePellStepDown.index_residues} combines congruences
that determine quantities only up to sign. The original proof introduces a
local sign parameter, splits two alternatives against two more alternatives,
extracts divisibility witnesses, and manufactures their sums or differences.

All this arithmetic expresses transitivity of the relation
\[
 a\sim_m b
 \quad\Longleftrightarrow\quad
 a\equiv b\pmod m\ \lor\ a\equiv-b\pmod m.
\]
It is useful to give that relation a small explicit interface:

\begin{Lean}
abbrev SignedModEq (m a b : ℤ) : Prop :=
  a ≡ b [ZMOD m] ∨ a ≡ -b [ZMOD m]
\end{Lean}

Symmetry follows by negating a congruence where necessary. Transitivity follows
from the four sign products:
\[
 (+)(+)=+,\quad (+)(-)=-,\quad (-)(+)=-,\quad (-)(-)=+.
\]
The helper proof performs those four cases once, using the existing
\code{Int.ModEq.trans} and \code{Int.ModEq.neg} operations. Client proofs do
not reopen integer quotient witnesses.

The new 43-line module provides symmetry, transitivity, conversion of
\(a\equiv\varepsilon b\) when \(\varepsilon\in\{1,-1\}\), a divisibility
characterization, and the absolute-value bridge
\(a\sim_m\operatorname{natAbs}(a)\).
It is an abbreviation of the existing disjunction, so the public conclusions
do not acquire a new wrapper type.

\subsection{The exact step-down composition}

The proof already has:
\[
 u\sim_c J,\qquad
 u\equiv(-1)^h s\pmod c,\qquad
 c\mid s-p\ \lor\ c\mid s+p.
\]
Convert the middle fact to \(u\sim_c s\), and the last fact to \(s\sim_c p\).
Then compose \(J\sim_c u\sim_c s\sim_c p\). The delivered client is:

\begin{Lean}
exact SignedModEq.iff_dvd_sub_or_add.mp
  (((SignedModEq.symm hu_J).trans
    (SignedModEq.of_sign_mul hu_cc (neg_one_pow_eq_or ℤ h))).trans
    (SignedModEq.iff_dvd_sub_or_add.mpr hsp'))
\end{Lean}

Two more clients consume the absolute-value bridge:
\code{NativePellStepDown.signed_index_residues} and
\code{SignedPellResidue.signed_residue_cases}. Their explicit positive/negative
cases become composition with \code{SignedModEq.natAbs}.

The modulus is an arbitrary integer. No positivity hypothesis is hidden:
for modulus zero, congruence becomes equality and the relation becomes
equality up to sign; negative moduli also remain valid. The proof uses the
unrestricted integer congruence API. Care is required in the divisibility
bridge because \code{Int.modEq_iff_dvd} uses \(m\mid b-a\). The helper explicitly
reverses the congruences before simplifying \(a-(-b)=a+b\).
See \ml{Data/Int/ModEq} for these API conventions.

\subsection{Charge the abstraction honestly}

Taken together, odd-index reuse, removal of the unused local fact, and the
signed-residue refactor change 350 lines in the two existing files into
292 lines there plus 43 lines of shared support: \textbf{335 lines total},
a net saving of 15. The signed relation by itself is primarily a structural
simplification; its support cost absorbs much of the client saving.

The reason to adopt it is that sign ambiguity becomes an explicit operation
with simple composition laws. The reason to defer it would be a preference
for the smallest possible local patch when no further signed-residue clients
are planned. The report does not count the parity saving a second time, and
does not describe all 58 removed client lines as a net reduction.

\section{Derive inverse pairing graphs by existential projection}
\label{sec:pair}

\subsection{One invertible graph has several coordinate views}

The forward pairing relation already supplies all the mathematics needed by
the inverse-coordinate representations. For natural pairing:
\[
\begin{aligned}
(\exists b,\ z=\operatorname{pair}(a,b))
  &\iff a=(\operatorname{unpair}(z))_1,\\
(\exists a,\ z=\operatorname{pair}(a,b))
  &\iff b=(\operatorname{unpair}(z))_2.
\end{aligned}
\]
The new helpers prove these facts directly using the existing inverse laws.
For example:

\begin{Lean}
theorem exists_pair_right_iff (z a : ℕ) :
    (∃ b, z = Nat.pair a b) ↔ a = (Nat.unpair z).1 := by
  constructor
  · rintro ⟨b, rfl⟩
    simp
  · intro h
    exact ⟨(Nat.unpair z).2, by rw [h, Nat.pair_unpair]⟩
\end{Lean}

At the quadratic-certificate layer, \code{exists_certificate_iff} already
relates the certificate to the forward pairing equation. Normalize the finite
witness function using \code{Fin.exists_fin_succ_pi}; consume that existing
equivalence; then apply the appropriate projection lemma.
The proof does not need to build the quotient/slack vector again.

This applies to seven constructors across \src{PairingRepresentation},
\src{QuadraticNatPairing}, and \src{QuadraticTreeCodeGraphs}:
the two ordinary inverse coordinates, the quadratic forward graph, its two
inverse coordinates, and the two tree-code coordinates. The tree case also
uses the existing predecessor equation; the definition of the zero-code
predecessor remains unchanged.

\subsection{An actual quadratic client}

The shortened left-coordinate constructor retains its system and witness
dimensions exactly:

\begin{Lean}
def leftRepresentation :
    Rep (Fin 2) (fun v => v 0 = (Nat.unpair (v 1)).1) :=
  (System.scopedRep (a := 2) (k := 4) (system 0 2)).ofIff (by
    intro v
    simp only [holds_iff]
    change (∃ w : Fin 4 → ℕ,
      Certificate (v 1) (v 0) (w 0) (w 1) (w 2) (w 3)) ↔ _
    simp [Fin.exists_fin_succ_pi, exists_certificate_iff,
      GraphRep.exists_pair_right_iff])
\end{Lean}

The seven declaration blocks shrink from 97 to 54 lines before support costs.
The two projection helpers and their placement are charged to the complete
module totals, which shrink from \textbf{396 to 369 lines}, a net saving of 27.

More importantly, all seven exact replacement constructors have passed
equalities of the form
\begin{Lean}
theorem leftRepresentation_short_eq :
    leftRepresentation_short = leftRepresentation := rfl
\end{Lean}
against the pinned original definitions. These equalities concern the full
\code{Rep} values, including their expression and witness-count data.
They are stronger than merely proving the same represented predicate.

This proposal is distinct from earlier witness-block transport and
partial-decoder injectivity suggestions. The mathematical reuse is projection
of the forward coding graph and elimination of its existing certificate.
The finite-function theorem is a library tool used to expose that structure;
see \ml{Data/Fin/Tuple/Basic}.

\section{Consume the Boolean certificate before constructing a conditional}
\label{sec:boolean}

\src{QuadraticBoolean} already proves
\code{exists_zero_certificate_iff}. Its zero-test and conditional
representations nevertheless unpack certificate soundness and completeness
and manually assemble witness functions.

For the zero test, normalize the one-coordinate witness function and apply
that equivalence directly. For the conditional, normalize its two witnesses
as a Boolean value \(b\) and a slack \(e\). The selection equation does not
depend on \(e\), so move the existential for \(e\) across that conjunct.
The certificate equivalence then determines
\[
 b=\begin{cases}1&n=0,\\0&n\ne0.\end{cases}
\]
Substituting this value into the selection equation yields the desired
conditional graph.

The final proof step is:
\begin{Lean}
by_cases hn : v 1 = 0 <;>
  simp [Fin.exists_fin_succ_pi, exists_and_right,
    exists_zero_certificate_iff, hn, sub_eq_zero]
\end{Lean}

The existing two branches remain, but they are now the actual semantic
branches of the conditional. There is no second manual proof of the
zero-certificate specification and no hand-built witness tuple.

The complete module shrinks from \textbf{136 to 126 lines}, with no new
production helper. Both actual replacement definitions compile, and both
complete-value equality checks pass by \code{rfl}.
Thus the zero system, conditional system, witness counts, and expression
data are preserved.

This is not the earlier scanner Boolean-interface proposal: it concerns
existential elimination in quadratic register representations, using a
specification already present in the same file. It is also not a new
Boolean encoding or a reduction in the number of Diophantine variables.
The reduction is in proof construction.

\section{Optional: recognize the physical permutation as a positive tail}
\label{sec:permutation}

This small proposal concerns \src{WangTiles942Permutation}.
Its existing zero and successor equations describe
\[
 \operatorname{physicalPosition}
 =\operatorname{Fin.cons}\bigl(0,\ i\mapsto
 1+\operatorname{strictRank}(\operatorname{nonStartWeight},i)\bigr).
\]
The zero head cannot equal a positive tail value, and the tail is injective
by cancellation and the already-proved strict-rank injectivity.
\code{Fin.cons_injective_iff} packages precisely those two facts.

The proposed proof is:
\begin{Lean}
theorem physicalPosition_injective :
    Function.Injective physicalPosition := by
  have hform : physicalPosition =
      Fin.cons 0
        (fun i : Fin 941 => 1 + strictRank nonStartWeight i) :=
    funext (Fin.cases physicalPosition_zero physicalPosition_succ)
  rw [hform]
  exact Fin.cons_injective_iff.mpr ⟨by simp,
    fun i j h => strictRank_injective nonStartWeight_injective
      (Nat.add_left_cancel h)⟩
\end{Lean}

The preceding \code{nonStartWeight_injective} can also use composition directly:
\begin{Lean}
theorem nonStartWeight_injective :
    Function.Injective nonStartWeight :=
  tile_weight_injective.comp (Fin.succ_injective 941)
\end{Lean}

The delivered file would shrink from 264 to 255 lines.
The general zero-extension theorem with these exact mathematical ingredients
has compiled successfully. The actual WangTiles942/K215 import chain was
\textbf{not rebuilt}, so the full repository client remains a proposal pending
its normal compilation gate. Its patch is in the optional directory and is
not included in the main patch or 106-line total.

Earlier reports recognize the repository's existing recovered-index and
strict-rank arguments. Those constructions are not claimed as new here.
Only the specific final injectivity packaging and composition replacement
are additional. The relevant library interface is in
\ml{Data/Fin/Tuple/Basic}.

\section{Validation and its limits}
\label{sec:validation}

\subsection{Pinned environment and compilation}

The verification used Lean 4.32.0, compiler commit
\code{8c9756b28d64dab099da31a4c09229a9e6a2ef35}.
All nine dependency package commit hashes were checked against the pinned
\code{lake-manifest.json}; mathlib is at
\code{81a5d257c8e410db227a6665ed08f64fea08e997}.
Dependency objects were reused from the corresponding cache, while
\textbf{63 repository modules were compiled from the freshly downloaded
pinned sources}. No previously built ProveIt objects from another snapshot
were used as the repository baseline.

All \textbf{ten main proposed modules} were then compiled in a complete local
import tree with the changed modules taking precedence. The finalizer pair,
the three pairing modules, and the signed-residue/step-down/bootstrap chain
were checked together in dependency order. The original working tree remained
unchanged. This checks the complete changed source files, rather than only
isolated snippets.

\begin{center}
\small
\begin{tabularx}{\linewidth}{@{}lX@{}}
\toprule
Check & Result and scope\\
\midrule
Rational integral-power lemma & Passed; preserves the exact signed rational root.\\
Finalizer modules & Both complete revised modules passed together.\\
Small Pell criterion & Passed, including the original first-unit shape and sharpness witness.\\
Full Pell modules & Revised signed residue, step-down, and bootstrap all passed.\\
Pairing constructors & Seven exact clients and seven full-value \code{rfl} equalities passed.\\
Boolean constructors & Two exact clients and two full-value \code{rfl} equalities passed.\\
Complete representation files & All four revised modules passed.\\
Signed congruence interface & Complete new module passed; unrestricted integer modulus.\\
Permutation schema & Generic theorem passed; full WangTiles942 client not compiled.\\
\bottomrule
\end{tabularx}
\end{center}

The included logs retain compiler diagnostics and source hashes. Intermediate
elaboration issues were corrected before the final successful compilations;
the files and records in the final package correspond to those successful
versions. Local elapsed times are recorded for reproducibility but are not
used as before/after performance measurements.

\subsection{What the representation equalities establish}

The \code{Rep} constructors contain both mathematical correctness fields and
computational data. Merely proving that a shortened constructor describes the
same relation would not establish that it has the same witness count or
literal expression.

The nine equalities compare each complete shortened constructor with its
original and close by \code{rfl}. Lean's proof irrelevance permits the
correctness proofs to differ, while the remaining data are definitionally
equal. This gives a direct check on the concern that proof refactoring might
silently change the Diophantine representation or size accounting.

Those checks run against the original snapshot before patch application.
After the patch, the newly shared helper names are already present, so the
original baseline-check files should not simply be rerun in the modified
namespace. The replay instructions make this order explicit.

\subsection{Axiom and interface review}

An integrated \code{#print axioms} audit checks the shared lemmas, the altered
proofs, and selected downstream endpoints for finalizer equivalence,
Pell classification, and representation construction.
The printed dependencies are drawn from the standard set
\[
 \{\code{propext},\ \code{Classical.choice},\ \code{Quot.sound}\}.
\]
No \code{sorryAx} or new custom axiom is introduced. The old and new
\code{firstUnit_fundamental} have the same printed axiom set.
This is an audit of these checked declarations, not a statement that every
possible repository theorem was audited.

Independent source review checks that all pre-existing declaration signatures
are retained, that the added imports are acyclic, and that the finalizer norm
and value definitions are unchanged. The new signed congruence is an
abbreviation of an existing proposition. The patch does not change any
published equation, register layout, witness dimension, or consistency
hypothesis.

\subsection{What has not been established}

The whole repository and all 574 Lean files were not rebuilt. The user should
run the repository's usual CI and endpoint checks when integrating the patch.
The optional permutation client remains uncompiled against its complete
dependency chain. No elaboration-time speedup, smaller kernel proof term,
or improved Diophantine variable count follows from the source-line numbers.

The mathematical reformulations are checked in ordinary Lean arithmetic.
This report makes no new claim about the equivalence between Diophantine
non-solvability and \(\operatorname{Con}(\mathrm{ZFC})\); it preserves the
interfaces of existing proofs used in that program.

\section{Adoption order and package guide}

\subsection{Recommended order}

The lowest-coupling first step is the finalizer pair. It gives the largest
net reduction, retains its imports, and replaces duplicated root-sign
bookkeeping by a direct integral-root interface.

Next adopt the pairing and Boolean changes. They consume local mathematical
specifications already present, and the nine full-value equality checks give
strong evidence that computational data are unchanged.

The odd-index reuse can be adopted as a small direct cleanup. The main patch
combines it with the signed-congruence module; review that abstraction on
maintainability grounds, with its complete support cost visible. Adopt the
small-ordinate Pell criterion if relying on the already-imported general
Pell classification is consistent with the preferred exposition.

The permutation proposal is last: its generic argument is checked and its
patch is small, but its full module must pass the normal project gate before
adoption.

\subsection{Archive contents}

\begin{center}
\small
\begin{tabularx}{\linewidth}{@{}lX@{}}
\toprule
Path & Purpose\\
\midrule
\code{report/} & This LaTeX source and its compiled PDF.\\
\code{patches/main.patch} & Combined checked changes, relative to the pinned repository.\\
\code{proposed/} & Complete replacement modules and the new signed-congruence module.\\
\code{checks/baseline/} & Original-snapshot constructor comparisons with nine \code{rfl} checks.\\
\code{checks/standalone/} & Reusable mathematical prototypes and the Pell sharpness example.\\
\code{checks/IntegratedAxiomAudit.lean} & Axiom queries for the combined modified modules.\\
\code{evidence/} & Source metrics, hashes, module-build records, and final compiler logs.\\
\code{audit/} & Prior-proposal exclusion map, archive/member hashes, and peer review.\\
\code{optional/} & Permutation patch and exact-client candidate, explicitly not fully compiled.\\
\code{scripts/} & Patch checking, replay, and archive-integrity verification helpers.\\
\bottomrule
\end{tabularx}
\end{center}

The root README contains the concrete commands. Before applying a patch,
verify the commit and run \code{git apply --check}. The baseline constructor
checks belong before application; the integrated axiom audit belongs after
building the modified modules. The replay helper requires an explicit
\code{--apply} argument before it changes the supplied checkout.

The proposed source files are included for inspection, but the patch is the
integration mechanism. Do not independently copy files and then apply the
same patch. The optional permutation patch does not overlap the main patch.

\section{Further research questions}

\subsection{Can rational witness descent become a small common API?}

The new lemma already handles every nonzero natural power, rather than only
squares. Other proofs that create rational witnesses and later prove their
powers integral may benefit from returning the original signed witness.
A useful follow-up is to survey actual consumers before creating a broader
abstraction. An integrally closed domain formulation is mathematically
natural, but could introduce heavier imports than the elementary denominator
proof. The decision should be based on repeated clients, not generality alone.

\subsection{How far does the small-ordinate Pell interface extend?}

The absolute threshold three is sharp. A more informative criterion could
use additional lower bounds on a hypothetical fundamental unit, or a
discriminant-dependent bound. The useful question is whether another
existing minimality proof can consume such a criterion and offset its support
cost. It would also be reasonable to keep both the elementary special proof
and the short general proof if their different mathematical explanations
serve different documentation needs.

\subsection{Where else is sign ambiguity a compositional relation?}

The new relation models the two units \(1\) and \(-1\). More complicated
root-of-unity ambiguity can be modeled by a finite group action, but this
would be a different and more expensive abstraction. First identify existing
proofs that already branch repeatedly on signs. Additional signed-residue
clients are the most plausible way to turn the current structural
improvement into a larger source reduction.

\subsection{Can certificate consumers systematically preserve data by construction?}

The pairing and Boolean examples suggest a focused review question:
does a representation constructor reopen soundness and completeness of a
certificate after an existential specification has already been proved?
Such a client may be shortened without changing its system, expression,
or witness count. Every candidate should retain a full-value equality check
where possible. This guards against confusing a proof simplification with a
different encoding that merely represents the same predicate.

\subsection{Can future simplification reports compare maintained proof cost?}

Large generated declarations can dominate raw line counts, while a small
semantic API can remove more conceptual duplication. Future evaluations
could record three separate quantities: complete maintained-source lines,
the number of independent mathematical arguments, and measured
elaboration/kernel-checking costs. This report establishes only the first
quantity numerically and explains the second qualitatively. It supplies
hashes and replay files so a controlled performance comparison can be run
without conflating different repository snapshots.

\appendix
\section{Novelty comparison at a glance}

The full exclusion map is intentionally kept in the machine-readable and
Markdown audit files. This table records the nearest existing ideas so that
the distinction of each proposal can be reviewed directly.

\small
\begin{longtable}{@{}p{0.24\linewidth}p{0.31\linewidth}p{0.38\linewidth}@{}}
\toprule
Proposal & Nearby prior material & Additional content here\\
\midrule
\endhead
Exact rational-root descent &
Cast cleanup and common scalar finalizer tails &
Keep the actual rational root through denominator-one descent; remove four
sign combinations and share the separated reconstruction theorem.\\[5pt]
Small-ordinate Pell criterion &
Pell cast bridges and named source-register facts &
A sharp general \(0<y\le3\) fundamental-unit criterion with a checked
counterexample at \(y=4\).\\[5pt]
Odd-index reuse &
Whole-disjunction cast simplifications &
Use the existing \code{odd_index_of_base_dvd} theorem to delete a duplicate
parity proof in another module.\\[5pt]
Signed congruence composition &
Ordinary modular/carry interfaces &
Symmetry and transitivity of congruence up to sign, plus a natural
absolute-value bridge.\\[5pt]
Pairing graph projections &
Witness-block transport, nine-coordinate packaging, partial-decoder injection &
Eliminate a coordinate of the forward pairing graph through its inverse
laws and existing quadratic certificate iff theorem.\\[5pt]
Boolean certificate elimination &
Scanner Boolean interfaces and tuple-layout suggestions &
Consume \code{exists_zero_certificate_iff} before solving the quadratic
conditional graph.\\[5pt]
Physical-position injection &
Existing recovered-index and strict-rank proofs &
Use \code{Fin.cons_injective_iff} for the final zero-head, positive-tail
map; full client remains optional and uncompiled.\\
\bottomrule
\end{longtable}
\normalsize

The incoming README points to an \code{ADOPTION.md} file that is absent in
the inspected commit. Its available placed report, README, placement note,
and artifacts were inspected; no content is attributed to the missing file.
Archive hashes in the supplied audit make the comparison reproducible.

\section{Primary source anchors}

All project links in this document are pinned to the inspected commit.
The principal targets are \src{NearSquareFinalizer}, \src{MarginFinalizer},
\src{NativePellBootstrap}, \src{NativePellStepDown}, \src{SignedPellResidue},
\src{PairingRepresentation}, \src{QuadraticNatPairing},
\src{QuadraticTreeCodeGraphs}, \src{QuadraticBoolean}, and
\src{WangTiles942Permutation}.

The primary external APIs used in the new arguments are in the pinned
mathlib modules \ml{NumberTheory/Pell}, \ml{Data/Rat/Defs},
\ml{Data/Rat/Lemmas}, \ml{Data/Int/ModEq}, and
\ml{Data/Fin/Tuple/Basic}. The Lean core denominator-power identity is checked
using the exact toolchain recorded in the validation evidence. No changing
web documentation is needed to replay the proofs.
\end{document}
